\documentclass[a4paper,11pt]{article}
\usepackage[utf8]{inputenc}
\usepackage[russian]{babel}
\usepackage[T1]{fontenc}
\usepackage{amssymb,amsmath,graphicx,indentfirst}
\usepackage{caption}
\usepackage{color}
\usepackage{listings}
\usepackage[unicode]{hyperref}

\setlength{\parskip}{1ex plus 0.5ex minus 0.2ex}
\captionsetup[figure]{labelformat=empty}
\captionsetup[figure]{justification=centering}
\lstset{keywordstyle=\color{blue}\bfseries}
\lstset{extendedchars=false, language=Caml, defaultdialect=[Objective]Caml}

\author{Олег Смирнов\\
\texttt{oleg.smirnov@gmail.com}}
\date{30 сентября 2011 г.}
\title{Введение в функциональное программирование -- Лекция 3. Функция как 
объект манипулирования. Абстракция вычисления}

\begin{document}
\maketitle
\tableofcontents
\newpage

\section*{Цель лекции}
\begin{itemize}
\item Абстракция вычисления
\item Функции высших порядков
\item Каррирование и разрезы
\end{itemize}

\section{Абстракция вычисления}
В конце прошлой лекции был рассмотрен пример функции, вычисляющей квадратный
корень методом Ньютона. Эта задача может быть разбита на подзадачи меньшего
объёма:

\begin{enumerate}
\item Получить начальное приближение.
\item Итеративно повторять шаги $2$ и $3$.
\item Вычислить новое приближение
\item Проверить, является ли результат достаточно хорошим. Если да -- вернуть
ответ
\end{enumerate}

В результате такой декомпозиции мы можем написать реализацию метода Ньютона,
абстрагированную от конкретной задачи -- от вычисления квадратного корня.
Вынося, например, вычисление нового приближения в отдельную подзадачу, мы 
получим универсальную программу, с помощью которой можно вычислить приближение
любой математической функции.

Пусть метод Ньютона принимаеть функцию вычисления нового приближения в качестве
аргумента. Такая функция, манипулирующая другими функциями, называется функцией
высшего порядка (higher-order function):

\begin{lstlisting}
  let newton first guess enough =
      let rec iter cur = 
          if enough cur then cur
          else iter (guess cur)
      iter first
\end{lstlisting}

Теперь для вычисления квадратного корня достаточно реализовать несколько
вспомогательных функций и передать их функции $newton$ в качестве аргументов:

\begin{lstlisting}
  let isGood x y = abs (x - y*y) < 0.000001
  let improve x y = ((x / y) + y) / 2.0
  let sqrt x = newton 1. (improve x) (isGood x)
\end{lstlisting}

Вообще говоря, поддержка функций высшего порядка в языке означает способность
языка манипулировать описаниями вычислений как полноценными значениями наравне
с другими структурами данных. Для языков без такой способности существует ряд
приёмов для имитации функций высшего порядка.

Функции высшего порядка иногда называют ``операторами'' или ``функционалами''
по аналогии с терминами функционального анализа.

\section{Функции в качестве аргументов}
Запишем функцию для суммирования значений $i$ для $i = a \ldots b$:
\begin{lstlisting}
  let rec sum a b = 
      if a > b then 0
      else a + sum (a + 1) b
\end{lstlisting}

А теперь напишем функцию для суммирования квадратов:
\begin{lstlisting}
  let rec sumSq a b = 
      if a > b then 0
      else a*a + sumSq (a + 1) b
\end{lstlisting}

Обе функции могут быть переписаны в виде итеративного процесса, но в этом
примере это не важно. Код функций похож, как будто он написан с использованием
общего шаблона. Присутствие такого шаблона является веским доводом в пользу
того, что здесь скрыта полезная абстракция. Действительно, в математической
нотации эта абстракция называется ``сигма-записью'':

\begin{equation*}
  sum = \sum_{i=a}^{b} i
\end{equation*}

\begin{equation*}
  sumSq = \sum_{i=a}^{b} i^2
\end{equation*}

Сила сигма-записи в том, что она выражает само понятие суммы вообще, а не
просто некоторые конкретные суммы. Выразим эту идею в виде функции:
\begin{lstlisting}
  let rec sumOf op a next b =
      if a > b then 0
      else (op a) + sumOf op (next a) next b
\end{lstlisting}

Здесь функция next означает ``приращение аргумента''. Её можно использовать,
например, при вычислении суммы ряда 
\begin{equation*}
  \frac{\pi}{4} = \frac{1}{3 \cdot 5} + \frac{1}{7 \cdot 9} + 
  \frac{1}{9 \cdot 11} + \ldots
\end{equation*}

Выразим sum и sumSq через sumOf:
\begin{lstlisting}
  let id a = a
  let inc a = a + 1
  let square a = a*a
  let sum a b = sumOf id a inc b
  let sumSq a b = sumOf square a inc b
\end{lstlisting}

Такая запись на порядок понятней. Теперь sumOf выражает идею суммирования, а 
sum и sumSq -- используют эту идею. Перепишем функцию sumOf в виде рекурсивного
процесса:
\begin{lstlisting}
  let sumOf op a next b =
      let rec sumIter a s =
          if a > b then s
          else sumIter (next a) (s + (op a))
      sumIter a 0
\end{lstlisting}

Функции sum и sumSq автоматически стали итеративными процессами! Абстракция
вычислений позволяет ``спрятать'' обобщённый вычислительный процесс и изменять
его реализацию по мере надобности. Кроме того, идею суммы можно ещё сильнее
обобщить, например, добавить предикат для ``отфильтровывания'' некоторых термов,
обобщить условие останова и т.д. 

Если вычислить сумму $\frac{1}{n}$ в интервале от 1 до 100000000, то получим
ответ 15.403683. Это отличается от правильного ответа 18.997896 на 3.5! Причина
такой потери точности -- разрядность типа float. Если сложить $10^6$ и $10^{-6}$
типа float, то результат будет $10^6$, т.к. тип float не обладает точностью в
12 разрядов. Если к $10^6$ прибавить $10^{12}$ раз по $10^{-6}$, то результат всё
равно будет $10^6$, хотя правильный ответ $2\cdot10^6$.

Способ суммирования тем оптимальнее, чем реже мы складываем большое число с 
маленьким. Поэтому на первом шаге нужно сложить два самых маленьких числа, затем
сложить оставшиеся числа и результат предыдущей операции и т.д. Это та же самая
задача -- суммирование большого списка чисел, просто список изменился. Алгоритм
можно эффективно реализовать в виде итеративного процесса с использованием
приоритетной очереди, затем алгоритм встроить в функцию sumOf. Возможен ряд
дополнительных оптимизаций, например переключение между ``стандартным
алгоритмом'' и ``алгоритмом с очередью'' в зависимости от значения аргументов.

\subsection*{Лямбда-функции}
До сих пор мы явно декларировали функции, передаваемые в sumOf в качестве
параметров. Однако это не обязательное действие. Функцию можно определить в
виде выражения, не связывая с ним никакое имя. В синтаксисе F\# для этого 
используется ключевое слово $fun$ и конструкция $fun~\langle arg
\rangle \rightarrow \langle body \rangle$. Например, функцию sum можно было
определить так:
\begin{lstlisting}
  let sum a b = sumOf (fun x -> x) a (fun x -> x + 1) b
\end{lstlisting}

Такая нотация называется определением лямбда-функции или анонимной функции --
функцией без имени.
\subsection*{Операции над списками}
Списком в F\# называется упорядоченная последовательность значений одного типа.
Например, список значений типа $int$ можно определить как:
\begin{lstlisting}
  let listInts = [1; 2; 3; 100]
\end{lstlisting}
Список значений типа $bool$:
\begin{lstlisting}
  let listBools = [false; true; true; true]
\end{lstlisting}
Пустой список обозначается как $[]$. Для списков определена операция $::$,
позволяющая присоединить элемент к ``голове'' списка:
\begin{lstlisting}
  0 :: [3; 2; 1]
  val it : int list = [0; 3; 2; 1]
\end{lstlisting}
Два списка можно ``склеить'' в новый список операцией $@$:
\begin{lstlisting}
  [true; true] @ [false; true; false]
  val it : bool list = [true; true; false; true; false]
\end{lstlisting}
Список в F\# неизменяем в том смысле, что в нём невозможно изменить элемент.

Для работы со списками удобно использовать сопоставление с образцом:
\begin{lstlisting}
  match someList with
  | head :: tail -> tail
  | [] -> []
\end{lstlisting}

Такая операция сначала сравнит список $someList$ с образцом $head :: tail$. 
Если это возможно, и список представим в виде конкатенации (склейки) головного 
элемента $head$ и списка из остальных элементов (хвоста) $tail$, то выполнится
первая ветка и результатом операции будет список-хвост $tail$. Если же нет, и
список $someList$ не содержит элементов (т.е. совпадает с $[]$), результатом
операции будет пустой список.

Напишем функцию суммирования списка в рекурсивном и итеративном варианте:
\begin{lstlisting}
  let rec sum list =
      match list with
      | head :: tail -> head + sum tail
      | [] -> 0  
\end{lstlisting}

\begin{lstlisting}
  let sum list =
      let rec loop list acc =
          match list with
          | head :: tail -> loop tail (acc + head)
          | [] -> acc
      loop list 0
\end{lstlisting}

Аналогично суммированию чисел в интервале, здесь можно выделить общую идею -- 
``свёртку'' списка с помощью некоторой бинарной операции $op$ и начального 
значения $base$. Функция свёртки применяет $op$ к двум аргументам -- к $base$ и
к первому элементу списка $x_1$. На следующем шаге $op$ применяется к результату
предыдущего действия и элементу $x_2$ и так далее, в результате вычисляя формулу:
\begin{equation*}
  op~(\ldots (op~(op~base~x_1)~x_2) \ldots)~x_n
\end{equation*}

В функциональном программировании эту операцию обычно называют $fold$. Различают
``левую'' и ``правую'' свертки ($foldl$ и $foldr$) или же ``свёртку вперёд'' и
``свёртку назад'' ($fold$ и $foldBack$). Разница в том, что левая или обратная
свёртка вычисляет формулу
\begin{equation*}
  op~x_1~(op~x_2~(\ldots (op~x_n~base) \ldots))
\end{equation*}

Результаты левой и правой свёртки будут совпадать, если операция $op$ обладает
свойством ассоциативности.

Пусть функция свёртки опредена как $fold~op~base~list$. Тогда на её основе можно
определить ряд полезных операций над списками. Например, суммирование элементов
списка:
\begin{lstlisting}
  let sumList list = fold (fun a b -> a + b) 0 list
\end{lstlisting}

Минимум и максимум в списке:
\begin{lstlisting}
  let min a b = 
      if a < b then a
      else b
  let max a b = 
      if a > b then a
      else b
  let minList list = fold min infinity list
  let maxList list = fold max (-infinity) list
\end{lstlisting}

Функции $all$ и $any$ для списков типа $bool$:
\begin{lstlisting}
  let all list = fold (fun a b -> a && b) true list
  let any list = fold (fun a b -> a || b) false list
\end{lstlisting}

\section{Функции в качестве результата}
Функции высшего порядка могут не только принимать другие функции в качестве
аргумента, но и возвращать их в качестве результата.

Определим несколько полезных операторов:
\begin{itemize}
\item Тождественный оператор:
  \begin{lstlisting}
    let id f x = f x
  \end{lstlisting}
\item Оператор композиции функций:
  \begin{lstlisting}
    let inline (<<) f g x = f (g x)
  \end{lstlisting}
\item Обратная композиция функций:
  \begin{lstlisting}
    let inline (>>) f g x = g (f x)
  \end{lstlisting}
\item Оператор применения функции:
  \begin{lstlisting}
    let inline ($) f x = f x
  \end{lstlisting}
\item Оператор смены порядка аргументов $flip$:
  \begin{lstlisting}
    let flip f x y = f y x
  \end{lstlisting}
\item Оператор вычисления в заданной точке $at$:
  \begin{lstlisting}
    let inline (|>) x f = f x
  \end{lstlisting}
  Другое определение:
  \begin{lstlisting}
    let at x f = (flip ($)) x f
  \end{lstlisting}
\end{itemize}

Ключевое слово $inline$ в определении композиции функций позволяет использовать
этот оператор в инфиксной форме, т.е. ставит его между аргументами. Оператор
обратной композиции часто бывает удобней, т.к. он применяет функции в том порядке,
в котором они записаны:
\begin{lstlisting}
  let incAndSquare x = (inc >> square) x
\end{lstlisting}

С помощью оператора вычисления в заданной точке удобно строить конвеерную
обработку, значительно упрощая вид функции. Сравним:
\begin{lstlisting}
  let complexFun x = square (add 5 (cube x))
  let complexFun x = x |> cube |> add 5 |> square
\end{lstlisting}

\begin{itemize}
\item Оператор отображения списка $map$ определяется следующим образом:
\begin{lstlisting}
  let rec map op list =
      match list with
      | head :: tail -> op head :: map op tail
      | [] -> []
\end{lstlisting}

\item Оператор фильтрации списка $filter$:
\begin{lstlisting}
  let rec filter pred list =
      match list with
      | head :: tail -> if pred head then
                           head :: filter pred tail
                        else filter pred tail
      | [] -> []
\end{lstlisting}
\end{itemize}

Теперь, комбинируя функции с помощью оператора композиции, легко построить
потоковую обработку списка. Например, функция, которая возводит каждый
элемент списка в четвёртую степень, предварительно отбросив отрицательные
элементы:
\begin{lstlisting}
  let filterNegAndQuad list = 
      list |> filter (fun x -> x > 0) |> map (square >> square)
\end{lstlisting}

Очевидно, что такой код гораздо более читаем, чем ``лесенка'' из вложенных
циклов в императивном стиле.

На практике функции высшего порядка имеют множество применений. Например, 
создание подпроцесса (нити, потока), реализация обобщённых алгоритмов
(например, сортировки, куда передаётся функция сравнения элементов) и
передача функции обратного вызова (callback) в событийно-ориентированном
программировании (event-driven programming).

\section{Каррирование и разрезы}
В предыдущем примере был построен ``конвеер'' из оператора фильтрации и
оператора отображения списка, причём в каждом из них первый аргумент был 
``зафиксирован''. Для фильтра был зафиксирован предикат 
$fun~x~\rightarrow~x > 0$, а для отображения -- функция возведения в
четвёртую степень $square~>>~square$.

Такой приём называется ``каррингом'' или частичным вычислением функции.
По определению, частичным применением $n$-арной функции $f$ называется 
конструкция, значением которой является $(n-k)$-арная функция, 
соответствующая $f$ при фиксированных значениях некоторых $k$ из $n$ её
аргументов. 

На практике имеет смысл упорядочивать аргументы функции по мере наростания
``изменчивости'': вначале идут ``аргументы-опции'', а затем ``аргументы-
данные''. Например, в описанном выше операторе $filter$ первым идёт
параметр-предикат $pred$, а затем список $list$. С точки зрения программиста,
удобно зафиксировать фильтрацию с предикатом в виде функции от одного 
аргумента, а затем применять её к различным спискам.

Некоторые языки позволяют превратить любую функцию от двух аргументов
в бинарный оператор. Фиксация левого или правого аргумента такого бинарного
оператора называется левым и правым разрезом функции соответственно.

\section{Бесточечный стиль}
В бесточечном стиле программирования мы опускаем аргументы (точки применения)
функции в левой и в правой частях ее определения. То есть вместо
\begin{lstlisting}
  sum list = fold add 0 list
\end{lstlisting}
в бесточечном стиле мы пишем
\begin{lstlisting}
  sum = fold add 0  
\end{lstlisting}

Такое отбрасывание аргументов возможно, только если самый правый аргумент в
левой части является самым правым и в правой части (и при этом больше нигде в
правой не используется). В нашем примере это аргумент $list$.

В бесточечном стиле функции не определяются через результат вычислений над 
аргументами, а выражаются через другие функции при помощи операторов композиции
функции и подобных. Такой стиль часто позволяет уменьшить количество лишней
информации в определении, делая его более лаконичным и читабельным.

Понятие карринга было впервые введено русским математиком Шейнфинкелем в
1924 г. при создании комбинаторной логики. Хаскелл Карри упоминал его в своих 
работах как ``первооткрывателя''. Бесточечный стиль в современном понимании
этого слова был впервые описан Джоном Бэкусом в его лекции ``Can programming
be liberated from the Von Neumann style?''~\cite{Backus} прочитанной на
церемонии вручения премии Тьюринга в 1977 году. 
\begin{figure}[h]
  \begin{minipage}[h]{0.49\linewidth}
    \begin{center}
        \includegraphics[height=40mm]{lecture3/schonfinkel.eps}
        \caption{Моисей Исаевич Шейнфинкель\\(1889 - 1942)}
    \end{center}
  \end{minipage}
  \begin{minipage}[h]{0.49\linewidth}
    \begin{center}
      \includegraphics[height=40mm]{lecture3/backus.eps}
      \caption{John Backus\\(1924 - 2007)}
    \end{center}
  \end{minipage}
\end{figure}

Обычно в языках программирования существуют ограничения на способы, с помощью
которых можно манипулировать элементами вычисления. Элементы или значения, на
которые накладывается наименьшее число ограничений, ещё называются
``полноправными'' или элементами первого класса (first-class). Вот некоторые из
их ``прав и  привилегий'':
\begin{itemize}
\item Их можно называть с помощью переменных.
\item Их можно передавать в процедуры в качестве аргументов.
\item Их можно возвращать из процедур в виде результата.
\item Их можно включать в структуры данных
\end{itemize}
\nocite{*}
\bibliographystyle{alpha}
\bibliography{lecture3}
\end{document}
